% appendix

\chapter{Feature Extraction C Source Code}
\label{app:feature_extraction}

The following C implementation of the feature extraction pipeline (Section~\ref{sec:ch3_features} of Chapter~\ref{methodology}) runs on the SPC58EC80E5 microcontroller. It receives arrays of raw measurements (voltage, current, temperature, time) for one window and produces the ten features listed in Table~\ref{tab:intermediate_signals}. The code avoids dynamic memory allocation (the MCU environment has no \texttt{malloc} support in this configuration); intermediate buffers are declared as stack arrays sized to the window length.

\begin{lstlisting}[language=C, basicstyle=\small\ttfamily, columns=fullflexible]
#include "compute_features.h"


void compute_features(float* voltage, float* current, float* temperature, float* time,
    int N, float* out_features)
{
    /*
    'mean_discharge_voltage_rate', 'mean_voltage',
    'mean_discharge_temperature_rate', 'mean_power',
    'mean_relativeTime', 'mean_dV_dQ',
    'max_dV_dQ', 'min_dV_dQ', 'duration',
    'mean_temperature'
    */

    float sum_discharge_voltage_rate = 0;
    float sum_voltage = voltage[0];
    float sum_discharge_temperature_rate = 0;
    float sum_power = voltage[0] * current[0];
    float sum_relativeTime = 0;
    float sum_dV_dQ = 0;
    float max_dV_dQ = -3.4028235e+38f;  // the min number that can be represented by float type
    float min_dV_dQ = 3.4028235e+38f;  // the max number that can be represented by float type
    float duration = 0;     // = max_relativetime - min_relativetime
    float sum_temperature = temperature[0];

    // set the reference of 'time' as 0
    float time_0 = time[0];
    for(int i=0; i<N; i++)
    {
        time[i] -= time_0;
    }

    // dQ
    // our device doesn't support 'malloc':
    // float* delta_current = (float*)malloc((N-1)*sizeof(float));
    float delta_current[N];
    float dQ[N];
    cumulative_trapezoid(time, current, delta_current, N);
    for(int i=0; i<N-1; i++)
    {
        delta_current[i] = delta_current[i] / 3600.0 * -1.0;
    }
    compute_gradient(delta_current, dQ, N-1);

    //dV
    float dV[N];
    compute_gradient(voltage, dV, N);


    for (int i=1; i<N; i++) {
        sum_discharge_voltage_rate += (voltage[i]-voltage[i-1]) / (time[i]-time[i-1]);
        sum_voltage += voltage[i];
        sum_discharge_temperature_rate += (temperature[i]-temperature[i-1]) / (time[i]-time[i-1]);
        sum_power += voltage[i] * current[i];
        sum_relativeTime += time[i];
        float dV_dQ = dV[i-1] / dQ[i-1];
        sum_dV_dQ += dV_dQ;
        if (dV_dQ > max_dV_dQ) {max_dV_dQ = dV_dQ;}
        if (dV_dQ < min_dV_dQ) {min_dV_dQ = dV_dQ;}
        sum_temperature += temperature[i];
    }

    float mean_discharge_voltage_rate = sum_discharge_voltage_rate / (N-1);
    float mean_voltage = sum_voltage / N;
    float mean_discharge_temperature_rate = sum_discharge_temperature_rate / (N-1);
    float mean_power = sum_power / N; //
    float mean_relativeTime = sum_relativeTime / N;
    float mean_dV_dQ = sum_dV_dQ / (N-1);
    duration = time[N-1];
    float mean_temperature = sum_temperature / N;

    out_features[0] = mean_discharge_voltage_rate;
    out_features[1] = mean_voltage;
    out_features[2] = mean_discharge_temperature_rate;
    out_features[3] = mean_power;
    out_features[4] = mean_relativeTime;
    out_features[5] = mean_dV_dQ;
    out_features[6] = max_dV_dQ;
    out_features[7] = min_dV_dQ;
    out_features[8] = duration;
    out_features[9] = mean_temperature;
}


void compute_gradient(float* x, float* gradient, int length)
{
    gradient[0] = x[1] - x[0];
    gradient[length-1] = x[length-1] - x[length-2];
    for(int i=1; i<length-1; i++)
    {
        gradient[i] = (x[i+1]-x[i-1])/2.0f;
    }
}

void cumulative_trapezoid(float* x, float* y, float* ct_array, int length)
{
    float ct = 0;
    for(int i=0; i<length-1; i++)
    {
        float t = (y[i]+y[i+1])*(x[i+1]-x[i])/2.0f;
        ct += t;
        ct_array[i] = ct;
    }
}
\end{lstlisting}


\chapter{Cycle-Count Benchmark Code}
\label{app:benchmark}

The CPU cycle counts reported in Chapter~\ref{experiment} are obtained with the MCU's system timer (STMD3), which runs at 90\,MHz---half the 180\,MHz core clock. The measured code region is wrapped between \texttt{benchmark\_start()} and \texttt{benchmark\_end()} calls; the returned tick count is doubled to obtain CPU cycles.

\begin{lstlisting}[language=C, basicstyle=\small\ttfamily, columns=fullflexible]
void benchmark_start()
{
  stm_lld_resetcounter(&STMD3);
}

uint32_t benchmark_end()
{
  uint32_t elapsed = stm_lld_getcounter(&STMD3);
  return elapsed;
}

benchmark_start();
compute_features(voltage, current, temperature, relativeTime, N_length, out_features);
fe_cycles = benchmark_end() * 2;

// --> predict capacity
benchmark_start();
float pred_capacity = ensemble_inference(CHILDREN_RIGHT, ALPHAS, FEATURES, ROOTS, out_features);
capacity_cycles = benchmark_end() * 2;
\end{lstlisting}


\chapter{Python Dependencies}
\label{app:dependencies}

The training and evaluation pipeline was developed in Python. The complete dependency list of the training environment, as reported by \texttt{pip freeze}, is given in Table~\ref{tab:dependencies}. The core packages are LightGBM 4.5.0 (model training), Optuna 4.2.0 (hyperparameter search), and TensorFlow 2.18.0 with Keras 3.8.0 (BiLSTM baseline); NumPy and Pandas provide data handling.

\begin{table}[H]
    \centering
    \setcellgapes{3pt}
    \makegapedcells
    \begin{tabular}{c c c c}
    \hline
    \multicolumn{4}{c}{Core packages} \\
    \hline
    lightgbm & 4.5.0 & optuna & 4.2.0 \\
    tensorflow & 2.18.0 & keras & 3.8.0 \\
    numpy & 2.0.2 & pandas & 2.2.3 \\
    scikit-learn & 1.6.1 & scipy & 1.15.1 \\
    matplotlib & 3.10.0 & seaborn & 0.13.2 \\
    \hline
    \end{tabular}
    \caption{Core Python packages of the training environment}
    \label{tab:dependencies}
\end{table}



